% Self-contained mathematical sections, included by report.tex.
\section{The retained system and exact scaling}
\label{sec:exact-full-system}
The analysis keeps the intracellular and luminal amounts of sodium,
potassium, chloride, total inorganic carbon and total alkalinity; both
volumes; and the AE4 regulatory fraction.  Carbonate speciation and two
membrane voltages are recovered algebraically at every evaluation.
The older reduction in \texttt{Marty.tex} is therefore not used: neither a
fixed pH nor a fixed volume can be inserted into the following identities.
Alkalinity is a conserved acid--base coordinate, not bicarbonate alone.

Let $N,H,E,B,J_4$ denote inward NKCC1 cycles, NHE1 sodium uptake with proton
extrusion, AE2 chloride uptake with bicarbonate extrusion, NBC cycles and
AE4 cycles respectively.  A positive NBC cycle imports one sodium and two
bicarbonates.  Write $P=P_a+P_b$ for the total pump cycles,
$K_a,K_b$ for potassium efflux from cell to lumen and bath, and $C_a$ for
chloride efflux from cell to lumen.  The paracellular flux $j_{p,s}$ is
positive from lumen to bath for species $s$.  The carbon dioxide fluxes
$D_b,D_a$ are positive from bath and lumen into the cell.  Water fluxes
$q_b,q_a,q_t,Q$ are positive from bath to cell, cell to lumen, bath to lumen,
and lumen to sink respectively.  These are signed fluxes, not their positive
parts.  The general AE4 source vector is
\[
 \nu=(\nu_{\mathrm{Na}},\nu_{\mathrm K},\nu_{\mathrm{Cl}},\nu_T,\nu_A),
 \qquad S^4=J_4\nu .
\]

Take $C_0=[\mathrm{Na}]_e$, $V_0=V_i(0)$, $L_0=V_l(0)$ and the initial
stimulated WT NKCC1 \emph{cycle} flux $J_0$.  Define
\[
 t_0=\frac{C_0V_0}{J_0},\quad \tau=t/t_0,\quad
 \varepsilon=L_0/V_0,\quad v=V_i/V_0,\quad \ell=V_l/L_0,
 \quad x_s=\frac{n_{s,i}}{C_0V_0},\quad y_s=\frac{n_{s,l}}{C_0L_0}.
\]
Every ionic or carbon flux below is divided by $J_0$, and every water flux
by $J_0/C_0$; a hat marks these fluxes.  Dimensionless concentrations are
$c_{s,i}=x_s/v$ and $c_{s,l}=y_s/\ell$.  Primes mean $d/d\tau$.
The full dimensionless amount and volume equations are
\begin{align}
 x_{\mathrm{Na}}'&=\widehat N+\widehat H+\widehat B+
                \nu_{\mathrm{Na}}\widehat J_4-3\widehat P,\notag\\
 x_{\mathrm K}'&=\widehat N+\nu_{\mathrm K}\widehat J_4+
                2\widehat P-\widehat K_a-\widehat K_b,\notag\\
 x_{\mathrm{Cl}}'&=2\widehat N+\widehat E+
                \nu_{\mathrm{Cl}}\widehat J_4-\widehat C_a,\notag\\
 x_T'&=-\widehat E+2\widehat B+\nu_T\widehat J_4+
                \widehat D_b+\widehat D_a,\notag\\
 x_A'&=\widehat H-\widehat E+2\widehat B+\nu_A\widehat J_4,\notag\\
 v'&=\widehat q_b-\widehat q_a,\label{eq:dimless-cell}\\
 \varepsilon y_{\mathrm{Na}}'&=3\widehat P_a-\widehat j_{p,\mathrm{Na}}
                -\widehat Q c_{\mathrm{Na},l},\notag\\
 \varepsilon y_{\mathrm K}'&=-2\widehat P_a+\widehat K_a-
                \widehat j_{p,\mathrm K}-\widehat Qc_{\mathrm K,l},\notag\\
 \varepsilon y_{\mathrm{Cl}}'&=\widehat C_a-\widehat j_{p,\mathrm{Cl}}
                -\widehat Qc_{\mathrm{Cl},l},\notag\\
 \varepsilon y_T'&=-\widehat j_{p,\mathrm{HCO}_3}-\widehat D_a-
                \widehat Qc_{T,l},\notag\\
 \varepsilon y_A'&=-\widehat j_{p,\mathrm{HCO}_3}-\widehat Qc_{A,l},\notag\\
 \varepsilon\ell'&=\widehat q_a+\widehat q_t-\widehat Q,
 \qquad \theta_4 r'=\beta-r,\qquad \theta_4=30/t_0.
 \label{eq:dimless-lumen}
\end{align}
The neutral manifold is
\begin{equation}
 x_A=x_{\mathrm{Na}}+x_{\mathrm K}-x_{\mathrm{Cl}}-a_f,
 \quad y_A=y_{\mathrm{Na}}+y_{\mathrm K}-y_{\mathrm{Cl}},
 \quad a_f=\frac{n_{\mathrm{fixed}}}{C_0V_0}.
 \label{eq:neutral-manifold}
\end{equation}
Eliminating these two amounts leaves eleven independent dynamic
coordinates.  This is exact: cell charge is conserved; luminal charge obeys
$\dot q_l=-Qq_l/V_l$ and hence remains zero from zero initial charge.
Alkalinity and its effect on pH remain in the reconstructed state.

\subsection{Dimensionless constitutive laws}
Here $h_i$ and $b_i$ denote the dimensionless hydrogen and bicarbonate
concentrations from the carbonate closure below.  All half concentrations
and dissociation constants in this subsection are divided by $C_0$;
numerical rate constants retain their declared unit conversions.
For constant full stimulation the frozen NKCC1 and NHE1 multipliers are
$m_N=1.75$ and $m_H=2.3$.  More generally their separate clipped calcium
coordinates have the frozen endpoints used by the production objects:
the active NKCC1 endpoints are $0.058$ and $0.1\,\mu\mathrm M$,
whereas NBC and NHE1 use $0.058$ and $0.25\,\mu\mathrm M$.
This distinction matters away from full stimulation.

Put $X=c_{\mathrm{Na},i}c_{\mathrm K,i}c_{\mathrm{Cl},i}^{2}$.  The
Palk law actually selected in Task 40~\cite{Palk2010} is
\begin{equation}
 \widehat N=e_Nm_N\frac{\alpha_N A_1}{J_0A_3}
   \frac{1-\eta_2X}{1+\eta_4X},\qquad
 \eta_2=\frac{A_2C_0^4}{A_1},\quad
 \eta_4=\frac{A_4C_0^4}{A_3}.
 \label{eq:dimless-nkcc}
\end{equation}
The coefficients are $A_1=157.5$, $A_2=2.0096\,10^{-5}$,
$A_3=1.0306$ and $A_4=1.3852\,10^{-6}$ in the repository's mM
convention.  The effective amount flux $\alpha_N$ comes from the frozen
WT calibration, not an imported membrane density.

For NHE1, choose $k_*=1000k_{1+}$, converting its printed rate per
millisecond into a rate per second.  Let barred rates be divided by $k_*$:
\begin{align*}
 \bar a&=\frac{1000k_{1+}}{k_*}
       \frac{c_{\mathrm{Na},e}}{c_{\mathrm{Na},e}+K_{\mathrm{Na},e}}
       \frac{K_{H,e}}{h_e+K_{H,e}},&
 \bar b&=\frac{1000k_{2+}}{k_*}
       \frac{h_i}{h_i+K_{H,i}}
       \frac{K_{\mathrm{Na},i}}{c_{\mathrm{Na},i}+K_{\mathrm{Na},i}},\\
 \bar c&=\frac{1000k_{1-}}{k_*}
       \frac{c_{\mathrm{Na},i}}{c_{\mathrm{Na},i}+K_{\mathrm{Na},i}}
       \frac{K_{H,i}}{h_i+K_{H,i}},&
 \bar d&=\frac{1000k_{2-}}{k_*}
       \frac{h_e}{h_e+K_{H,e}}
       \frac{K_{\mathrm{Na},e}}{c_{\mathrm{Na},e}+K_{\mathrm{Na},e}}.
\end{align*}
Thus, with the exact printed Cha constants selected by production~\cite{Cha2009},
\[
 M_H=\frac{1}{1+(1+h_e/K_{M,e})(K_{M,i}/h_i)^3},\quad
 \widehat H=\frac{n_Hk_*}{J_0}e_Hm_HM_H
        \frac{\bar a\bar b-\bar c\bar d}{\bar a+\bar b+\bar c+\bar d}.
\]
The printed $k_{2-}=183\,\mathrm{ms}^{-1}$ is retained: the small
rounding discrepancy in microscopic reversibility is not silently repaired.
The other acid--base flux laws are
\begin{align*}
 \widehat E&=e_2\gamma_2\tanh\left[
  \frac{\log(c_{\mathrm{Cl},e}b_i/(c_{\mathrm{Cl},i}b_e))}{w_2}\right],
 &\gamma_2&=G_2/J_0,\\
 \widehat B&=\gamma_Bu\tanh\left[
  \frac{\log(c_{\mathrm{Na},e}b_e^2/(c_{\mathrm{Na},i}b_i^2))+\psi_b}{w_B}\right],
 &\gamma_B&=G_B/J_0,\\
 \widehat D_b&=\delta_b(c_{\mathrm{CO}_2,e}-c_{\mathrm{CO}_2,i}),
 &\widehat D_a&=\delta_a(c_{\mathrm{CO}_2,l}-c_{\mathrm{CO}_2,i}).
\end{align*}
Here $\delta_{a,b}=P_{\mathrm{CO}_2,a,b}C_0/J_0$,
$\psi_b=V_b/V_T$, $V_T=RT/F$, and $u$ is the frozen NBC calcium
coordinate.  The tanh widths are dimensionless effective parameters.

For the inherited AE4 shared carrier, write $k_4$ for its common chloride
attempt rate.  Define
\begin{align*}
 L_c&=\log(c_{\mathrm{Cl},e}/c_{\mathrm{Cl},i}),&
 L_s&=\log(c_{s,i}/c_{s,e})+2\log(b_i/b_e),\quad s\in\{\mathrm{Na},\mathrm K\},\\
 a_c^\pm&=\exp(\pm L_c/2),&
 a_s^\pm&=(k_s/k_4)\exp(\pm L_s/2),\\
 U&=a_c^++a_{\mathrm{Na}}^-+a_{\mathrm K}^-,&
 W&=a_c^-+a_{\mathrm{Na}}^++a_{\mathrm K}^+,\quad
 \pi_i=U/(U+W),\quad\pi_o=W/(U+W).
\end{align*}
The scalar cycle is exactly
\[
 \widehat J_4=e_4\gamma_4(1+0.25r)
 \left[(a_{\mathrm{Na}}^++a_{\mathrm K}^+)\pi_i
       -(a_{\mathrm{Na}}^-+a_{\mathrm K}^-)\pi_o\right],
 \qquad\gamma_4=n_4k_4/J_0.
\]
The frozen cooperative gate is absent.  Task 40 applies equal ensemble
sodium and potassium routing to this scalar cycle.  The seven Catal\'an
classes change its source vector but do not provide new class specific
kinetics.  Local detailed balance of the inherited carrier edges does not
by itself establish thermodynamic consistency of those substituted vectors.

To specify the electrical laws, scale each conductance $G$ as
$g=GV_T10^{15}/(FJ_0)$.  The reversal potential for a valence $z_s$ is
$e_{s,uv}=z_s^{-1}\log(c_{s,v}/c_{s,u})$.  With
$g_{K,a},g_{K,b},g_{C,a}$ incorporating the frozen calcium gate
$a_{\mathrm{Ca}}$ and background conductances, explicitly
\[
 g_{K,a}=g_{K,\mathrm{total}}a_{\mathrm{Ca}}f_K+g_{a,\mathrm{bg}},\quad
 g_{K,b}=g_{K,\mathrm{total}}a_{\mathrm{Ca}}(1-f_K)+g_{b,\mathrm{bg}},\quad
 g_{C,a}=g_{\mathrm{Cl},a}a_{\mathrm{Ca}},
\]
the channel and paracellular fluxes are
\begin{align*}
 \widehat K_a&=g_{K,a}(\psi_a-e_{K,il}), &
 \widehat K_b&=g_{K,b}(\psi_b-e_{K,ie}),\\
 \widehat C_a&=-g_{C,a}(\psi_a-e_{\mathrm{Cl},il}), &
 \widehat j_{p,s}&=\frac{g_{p,s}}{z_s}
   (\psi_b-\psi_a-e_{s,le}).
\end{align*}
The apical pump is
\[
 \widehat P_a=\gamma_P f_P
 \frac{c_{\mathrm{Na},i}^3}{c_{\mathrm{Na},i}^3+K_{P,\mathrm{Na}}^3}
 \frac{c_{\mathrm K,l}^2}{c_{\mathrm K,l}^2+K_{P,\mathrm K}^2},
 \qquad\gamma_P=G_P/J_0,
\]
and the basolateral expression replaces $f_P,c_{\mathrm K,l}$ by
$1-f_P,c_{\mathrm K,e}$.  These pump rates are voltage independent in the
frozen model.  The calcium gate is $c_{\mathrm{Ca}}^n/(c_{\mathrm{Ca}}^n+K_{\mathrm{Ca}}^n)$,
using the same calcium units in numerator and denominator.

For water define dimensionless osmolarities
\[
 o_i=c_{\mathrm{Na},i}+c_{\mathrm K,i}+c_{\mathrm{Cl},i}+c_{T,i}
       +(\beta_i+\iota_i)/v,\quad
 o_l=c_{\mathrm{Na},l}+c_{\mathrm K,l}+c_{\mathrm{Cl},l}+c_{T,l},
\]
where $\beta_i=n_{\mathrm{buffer},i}/(C_0V_0)$ and
$\iota_i=n_{\mathrm{other\ impermeant}}/(C_0V_0)$.  The bath expression
includes its declared untracked osmolyte concentration.  Exactly as in the
production law, free hydrogen and hydroxide are omitted from osmolarity,
and the tiny luminal buffer is not added as an osmotic particle.  No new
approximation is introduced by the present scaling.  Then
\[
 \widehat q_b=\lambda_b(o_i-o_e),\quad
 \widehat q_a=\lambda_a(o_l-o_i),\quad
 \widehat q_t=\lambda_t(o_l-o_e),\quad
 \widehat Q=\omega\varepsilon(\ell-d_l)_+,
\]
with $\lambda_j=L_jC_0^2/J_0$, $\omega=k_{\mathrm{out}}t_0$,
and $d_l=V_{\mathrm{dead}}/L_0$.  These equations close
\eqref{eq:dimless-cell}--\eqref{eq:dimless-lumen} without fixing a volume,
composition, carbon pool or voltage.

\section{Carbonate closure and pH derivatives}
\label{sec:carbonate-proof}
Write $p$ for pH, $h=10^{-p}$ in molar units,
$K_1=10^{-pK_1}$ and $K_2=10^{-pK_2}$.  The carbonate fractions are
\[
 (\alpha_0,\alpha_1,\alpha_2)=
 \frac{(h^2,K_1h,K_1K_2)}{h^2+K_1h+K_1K_2},\qquad
 f(p)=\alpha_1+2\alpha_2,\quad
 g(p)=\frac1{1+10^{pK_b-p}}.
\]
In concentration units mM let
$w(p)=10^3(10^{p-pK_w}-10^{-p})=[\mathrm{OH}]-[\mathrm H]$.
For either compartment, with fixed finite buffer amount $n_b$, the closure
is the single equation
\begin{equation}
 G(p;n_T,n_A,V)=n_Tf(p)+n_bg(p)+Vw(p)-n_A=0.
 \label{eq:carbonate-amount}
\end{equation}
Equivalently, for the cell,
$x_A=x_Tf(p)+\beta_i g(p)+v w(p)/C_0$; the lumen uses $y_T,y_A,\ell$
and $\beta_l=n_{b,l}/(C_0L_0)$.  Thus dilution of the finite buffer is
retained exactly.

\paragraph{Proposition: unique speciation on the stated domain.}
Suppose $V>0$, $n_T\geq0$, $n_b\geq0$, and all constants are finite and
strictly positive in their concentration representation.  For each real
$n_A$, equation \eqref{eq:carbonate-amount} has exactly one real pH.
The solution depends smoothly on the retained amounts and volume on their
open physical domain.  For the production implementation one additionally
requires $n_T>0$ and
\[
 G(p_{\min})\leq0\leq G(p_{\max});
\]
smoothness within the executable domain is asserted with strict bracket
inequalities.  This numerical bracket is distinct from a physiological
pH acceptance interval.

\paragraph{Proof.}
Let $Z$ take the values $0,1,2$ with probabilities
$\alpha_0,\alpha_1,\alpha_2$.  Differentiation gives
\[
 f'(p)=\log(10)\operatorname{Var}(Z),\quad
 g'(p)=\log(10)g(1-g),\quad
 w'(p)=\log(10)([\mathrm H]+[\mathrm{OH}]).
\]
Consequently
\begin{equation}
 G_p=\log(10)\{n_T\operatorname{Var}(Z)+n_bg(1-g)
                  +V([\mathrm H]+[\mathrm{OH}])\}>0.
 \label{eq:carbonate-positive}
\end{equation}
The water term sends $G$ to $-\infty$ as $p\to-\infty$ and to
$+\infty$ as $p\to+\infty$.  Continuity, strict monotonicity and the
implicit function theorem prove the claims.  The statement for the
finite numerical bracket follows by monotonicity. $\square$

At fixed buffer amount the partial derivatives and trajectory derivative
are
\begin{equation}
 \frac{\partial p}{\partial n_A}=\frac1{G_p},\quad
 \frac{\partial p}{\partial n_T}=-\frac f{G_p},\quad
 \frac{\partial p}{\partial V}=-\frac w{G_p},\qquad
 \dot p=\frac{\dot n_A-f\dot n_T-w\dot V}{G_p}.
 \label{eq:ph-derivatives}
\end{equation}
If the buffer amount is a parameter, $\partial p/\partial n_b=-g/G_p$.
The very small explicit volume derivative at fixed amounts does not permit
constant volume: volume changes all transported concentrations, osmolarity,
electrical reversals and therefore $\dot n_T$ and $\dot n_A$.  Nor does
small free hydrogen concentration justify dropping carbon or buffering.

\section{The coupled electrical closure}
\label{sec:electrical-proof}
Use the dimensionless conductances and voltages defined above.  Set
$g_a=g_{K,a}+g_{C,a}$, $g_p=\sum_sg_{p,s}$,
$c_a=-g_{K,a}e_{K,il}-g_{C,a}e_{\mathrm{Cl},il}+\widehat P_a$
and $c_p=-\sum_sg_{p,s}e_{s,le}$.  The apical equation gives
\[
 \psi_a(\psi_b)=\frac{g_p\psi_b-c_a+c_p}{g_a+g_p}.
\]
The remaining current equation, including the electrogenic NBC, is
\begin{equation}
 R(\psi_b)=g_{K,b}(\psi_b-e_{K,ie})+\widehat P_b
      +\widehat B(\psi_b)+g_p(\psi_b-\psi_a(\psi_b))+c_p=0.
 \label{eq:electrical-scalar}
\end{equation}

\paragraph{Proposition: unique electrical root.}
At any chemical state with positive concentrations and finite logarithms,
suppose all conductances and NBC capacity are nonnegative,
$g_a+g_p>0$, and
$s_0=g_{K,b}+g_pg_a/(g_a+g_p)>0$.  If $w_B>0$ and $u\in[0,1]$,
the mathematical current closure has a unique root on the real voltage
line, smoothly dependent on the chemical state and parameters.  The
production routine further requires its physical voltage bracket
$[-0.5,0.5]\,\mathrm V$ to contain that root.

\paragraph{Proof.}
Writing $L_B=\log(c_{\mathrm{Na},e}b_e^2/(c_{\mathrm{Na},i}b_i^2))$,
\[
 R'(\psi_b)=s_0+\frac{\gamma_Bu}{w_B}
       \operatorname{sech}^2((L_B+\psi_b)/w_B)>0.
\]
The NBC term is bounded whereas the passive contribution grows linearly
with positive slope $s_0$.  Thus $R(\psi_b)\to\pm\infty$ at the two
ends of the real line.  Monotonicity and the implicit function theorem
complete the argument.  For any retained variable $z$,
$\partial_z\psi_b=-R_z/R_{\psi_b}$; differentiating the expression
for $\psi_a$ gives its corresponding derivative. $\square$

On a domain with positive amounts and volumes, strict pH and voltage
brackets and $V_l>V_{\mathrm{dead}}$, these closures make the constant input
chemical vector field smooth.  At the outflow threshold it is locally
Lipschitz but not differentiable.  The regulatory law on $0\leq r\leq1$
has the smooth analytical extension $\dot r=(\beta-r)/30$ and
capacity multiplier $1+0.25r$; its derivative at $r=1$ means the derivative
from the physical side, or equivalently that extension.  The executable
trial clamp outside $[0,1]$ must not be differentiated by a centred
difference across the boundary.  Doing so gives a spurious regulator
eigenvalue $-1/60$ instead of $-1/30\,\mathrm{s}^{-1}$.
These regularity statements prove local existence and uniqueness of
trajectories in the stated domain.  They do not prove global positivity,
global equilibrium uniqueness or global stability.

\section{Alkalinity balance and conditional capacity obstruction}
\label{sec:alkalinity-proof}
\paragraph{Proposition: parent alkalinity budget.}
For the parent equal routing class C0, the intracellular balance is
\begin{equation}
 \dot n_{A,i}=H+2B-E-2J_4.
 \label{eq:alkalinity-budget}
\end{equation}
Hence an intracellular stationary state must satisfy
$2J_4=H+2B-E$.  For $J_4\ne0$ define
\[
 \mathcal A=\frac{H+2B-E}{2J_4}
       =1+\frac{\dot n_{A,i}}{2J_4},\qquad
 \mathcal A_0=\frac{H-E}{2J_4}.
\]
Here $\mathcal A_0$ removes the NBC contribution at the same state; it is
not the result of resolving the NBC null system.  For $J_4=0$ these ratios
are undefined, and the relevant statement is the unnormalised balance
$\dot n_{A,i}=H+2B-E$.

\paragraph{Proof and finite time form.}
NHE1 proton extrusion adds one alkalinity equivalent, AE2 bicarbonate
export removes one, NBC imports two, and the parent AE4 cycle removes two.
Neutral carbon dioxide exchange, water transfer and the included potassium,
sodium and chloride pathways have no direct alkalinity source.
Integration yields, for every finite valid interval $[0,T]$,
\begin{equation}
 2\int_0^T J_4\,dt=\int_0^T(H+2B-E)\,dt
                         +n_{A,i}(0)-n_{A,i}(T).
 \label{eq:alkalinity-integrated}
\end{equation}
Changing volume is already retained because these are amount balances.
Equation \eqref{eq:ph-derivatives} then gives the pH change without
identifying alkalinity change with pH change. $\square$

\paragraph{Corollary: a conditional obstruction without NBC.}
Suppose a declared admissible domain gives $H\leq H_{\max}$ and
$E\geq-E_{\max}$.  At an equilibrium with $B=0$,
$J_4\leq(H_{\max}+E_{\max})/2$.  More generally, maintaining a positive
mean demand $J_d$ for time $T$ requires
\[
 2J_d\leq H_{\max}+E_{\max}
      +\frac{n_{A,i}(0)-n_{A,i}(T)}{T}
\]
when NBC is absent.  If $n_{A,i}(T)\geq A_{\min}$, a demand greater than
the capacity term can persist for at most
$[n_{A,i}(0)-A_{\min}]/[2J_d-H_{\max}-E_{\max}]$.
This is conditional on the stated domain and demand, not a universal
impossibility of secretion without NBC.

For the frozen fixed bath and NHE1 amount one can obtain a rigorous but
conservative bound from the inherited pH condition $p_i\geq6.6$.
In the dimensional Cha notation $a,d$ are fixed bath quantities,
$b\leq b_0(h)=1000k_{2+}h/(h+K_{H,i})$ and $c\geq0$.
Both $M_H(h)$ and $ab_0/(a+b_0+d)$ increase with $h$.
Therefore
\[
 H\leq n_Hm_H M_H(h_{\max})
      \frac{ab_0(h_{\max})}{a+b_0(h_{\max})+d},
 \qquad h_{\max}=10^3 10^{-6.6}\ \mathrm{mM}.
\]
With $m_H\leq2.3$ this yields $H_{\max}=0.08202591\,\mathrm{fmol/s}$;
the tanh AE2 law gives $E_{\max}=0.005\,\mathrm{fmol/s}$.
Consequently the corresponding NBC absent bound is
$J_4\leq0.04351295\,\mathrm{fmol/s}$.  A \emph{specified} demand of
$0.075\,\mathrm{fmol/s}$ requires
$B\geq0.03148705\,\mathrm{fmol/s}$ at equilibrium.
The chosen demand is an illustrative model loading level, not a measured
AE4 requirement.  Neither this bound nor the source based kinetic law
identifies the effective salivary carrier amounts.

\section{Equal routing and general source projections}
\label{sec:source-proof}
\paragraph{Proposition: exact chloride identity.}
For the full system with a constant AE4 source vector define
\[
 \kappa(\nu)=\nu_{\mathrm{Cl}}-2\nu_{\mathrm{Na}}+\nu_A.
\]
At every valid state,
\begin{equation}
 C_a=6P-H+\kappa(\nu)J_4
            +2\dot n_{\mathrm{Na},i}-\dot n_{A,i}-\dot n_{\mathrm{Cl},i}.
 \label{eq:general-chloride-identity}
\end{equation}
For the parent equal routing vector $\kappa=0$, yielding precisely
\[
 C_a=6P-H+2\dot n_{\mathrm{Na},i}-\dot n_{A,i}-\dot n_{\mathrm{Cl},i},
 \qquad C_a^*=6P^*-H^*.
\]

\paragraph{Proof.}
From the three intracellular balances,
\begin{align*}
 2\dot n_{\mathrm{Na},i}-\dot n_{A,i}-\dot n_{\mathrm{Cl},i}
 &=2(N+H+B+\nu_{\mathrm{Na}}J_4-3P)\\
 &\quad -(H-E+2B+\nu_AJ_4)
             -(2N+E+\nu_{\mathrm{Cl}}J_4-C_a)\\
 &=H-6P-\kappa(\nu)J_4+C_a.
\end{align*}
Rearrangement gives the claim.  Neither carbon, volume nor pH was frozen.
$\square$

Only the explicit AE4 term disappears from this stationary balance.
AE4 still changes the chemical state, pump, NHE1, NKCC1, pH, volumes,
membrane voltages and hence secretion.  The proposition does not prove
zero AE4 sensitivity, a small null phenotype for every admissible
parameterisation, or observational equivalence of mechanisms.  Away from
equilibrium the three storage terms cannot be omitted.  Integrating
\eqref{eq:general-chloride-identity} gives the analogous finite time
chloride budget by replacing each derivative integral with its endpoint
amount difference.

The five relevant linear projections of a source vector are
\[
 \Pi_{\mathrm{Cl}}\nu=\nu_{\mathrm{Cl}},\quad
 \Pi_T\nu=\nu_T,\quad \Pi_A\nu=\nu_A,\quad
 q\nu=\nu_{\mathrm{Na}}+\nu_{\mathrm K}-\nu_{\mathrm{Cl}}-\nu_A,
 \quad \kappa\nu=-2\nu_{\mathrm{Na}}+\nu_{\mathrm{Cl}}+\nu_A.
\]
The charge source is $(q\nu)J_4$.  Thus a class with nonzero $q\nu$ would
require its additional transported current in the electrical closure;
merely changing amount sources would not preserve the neutral manifold.
All seven frozen classes have $q\nu=0$, so their explicit AE4 current is
zero.  On that neutral subspace,
$\kappa\nu=\nu_{\mathrm K}-\nu_{\mathrm{Na}}$: the cancellation tests
equal sodium and potassium source coefficients.

\begin{center}
\begin{tabular}{lrrrrrr}
\hline
Class & $\nu_{\mathrm{Na}}$ & $\nu_{\mathrm K}$ & $\nu_{\mathrm{Cl}}$
      & $\nu_T$ & $\nu_A$ & $\kappa$\\
\hline
C0  & $-1/2$ & $-1/2$ & 1 & $-2$ & $-2$ & 0\\
C1  & 1 & $-1$ & 1 & $-1$ & $-1$ & $-2$\\
C2  & 1 & $-2$ & 1 & $-1$ & $-2$ & $-3$\\
C3a & 0 & $-1$ & 1 & $-2$ & $-2$ & $-1$\\
C3b & $-1$ & 0 & 1 & $-2$ & $-2$ & 1\\
C4a & 0 & $-1$ & 1 & $-1$ & $-2$ & $-1$\\
C4b & $-1$ & 0 & 1 & $-1$ & $-2$ & 1\\
\hline
\end{tabular}
\end{center}
Each individual balance annihilates the kernel of its displayed linear
functional.  For example a change in $\nu_T$ alone is invisible to
$\Pi_{\mathrm{Cl}},\Pi_A,q,\kappa$ but not to carbon balance.  This is
exactly why C3a/C4a and C3b/C4b have the same listed stationary chloride
coefficient but different carbon signatures.  Jointly the five projections
have rank five: if chloride, carbon and alkalinity projections vanish,
$q=0$ gives $\nu_{\mathrm K}=-\nu_{\mathrm{Na}}$, and $\kappa=0$
then gives $\nu_{\mathrm{Na}}=\nu_{\mathrm K}=0$.  No nonzero direction
is lost by all five together.  Equality in one projected balance is not
equality of molecular mechanisms or of the whole dynamical response.

\subsection{Independent verification and scaling implications}
The analysis script \texttt{verify\_exact.py} independently assembles
carbonate chemistry, all selected kinetic laws, the nonlinear electrical
solve, finite lumen transport, changing volumes and the thirteen amount,
volume and regulatory derivatives.  It compares these with production
at 84 random valid states spanning all seven source classes and varied
AE4 expression.  It also verifies the pH derivative formula at two
perturbation sizes, the positive electrical derivative, the exact charge
and chloride identities, and the correct inward regulatory derivative at
$r=1$.  Rational source coefficients are retained in
\texttt{output/exact\_verification.json}.  These implementation checks
support the algebra but do not turn random samples into a global theorem.

Numerically $\varepsilon=0.072515$ and $\theta_4=0.031977$;
$\gamma_B=0.515137$, $\gamma_2=0.022260$ and
$\gamma_4=0.668363$.  The water coefficients are
$\lambda_a=404.37$, $\lambda_b=4820.57$ and $\lambda_t=24.34$;
$\omega=187.64$.  These suggest some rapidly adjusting variables but
are insufficient to justify freezing either volume: their driving
osmotic differences and coupling to amounts matter.  The verified
dimensional recovery retains them all.
The small dissociation concentrations, luminal buffer amount and inherited
carrier relaxation ratio are recorded separately in the machine readable
groups.  A small buffer amount is a numerical convention here, not a
measured physiological asymptotic limit.  There is no membrane capacitance
or differential voltage equation in the frozen system, so no electrical
relaxation parameter can be inferred from its conductance groups.  The
algebraic voltage law is an inherited modelling assumption, not a new
singular perturbation theorem proved by this analysis.
